\documentclass[12pt,a4paper]{report}

% Main LaTeX file recreated from the supplied Stage-1 project report.
% Compile with pdfLaTeX in Overleaf.
% Put your figures in the same project folder and update filenames below.

\usepackage[T1]{fontenc}
\usepackage{mathptmx}
\usepackage[a4paper,margin=1in]{geometry}
\usepackage{graphicx}
\usepackage{tikz}
\usepackage{float}
\usepackage{array}
\usepackage{booktabs}
\usepackage{longtable}
\usepackage{ragged2e}
\usepackage{setspace}
\usepackage{titlesec}
\usepackage{enumitem}
\usepackage[hidelinks]{hyperref}
\usepackage{caption}

\onehalfspacing
\setlength{\parindent}{0pt}
\setlength{\parskip}{6pt}
\renewcommand{\arraystretch}{1.2}
\renewcommand{\contentsname}{Contents}
\renewcommand{\listfigurename}{List of Figures}

% ---------- EDITABLE PROJECT DETAILS ----------
\newcommand{\ProjectTitle}{INTELLIGENT PRE-TOLL PROCESSING AND DYNAMIC LANE MANAGEMENT SYSTEM}
\newcommand{\StudentOne}{Prathamesh Warkhede}
\newcommand{\RollOne}{404A041}
\newcommand{\StudentTwo}{Omkar Sonawane}
\newcommand{\RollTwo}{404A040}
\newcommand{\StudentThree}{Omkar Walse}
\newcommand{\RollThree}{404A039}
\newcommand{\GuideName}{Prof. R. R. Katkar}
\newcommand{\CoordinatorName}{Prof. V. B. Baru}
\newcommand{\HODName}{Dr. R. S. Kawitkar}
\newcommand{\PrincipalName}{Dr. S. D. Lokhande}
\newcommand{\CollegeName}{Sinhgad College of Engineering}
\newcommand{\CollegeAddress}{Vadgaon (Bk.), Pune -- 411041}
\newcommand{\AcademicYear}{2026--2027}

% ---------- OPTIONAL IMAGE HELPER ----------
% Displays an image if the file exists; otherwise prints a visible placeholder.
\newcommand{\projectimage}[2]{%
\IfFileExists{#1}{\includegraphics[width=#2\textwidth]{#1}}{%
\fbox{\parbox[c][4cm][c]{#2\textwidth}{\centering Image placeholder\\[2mm]\texttt{#1}}}}}

\begin{document}

% =========================================================
% COVER PAGE
% =========================================================
\begin{titlepage}
\thispagestyle{empty}
\begin{tikzpicture}[remember picture,overlay]
\draw[black!70,line width=3pt]
([xshift=1.2cm,yshift=-1.2cm]current page.north west)
rectangle
([xshift=-1.2cm,yshift=1.2cm]current page.south east);
\draw[black!70,line width=0.5pt]
([xshift=1.35cm,yshift=-1.35cm]current page.north west)
rectangle
([xshift=-1.35cm,yshift=1.35cm]current page.south east);
\end{tikzpicture}
\begin{center}
\vspace*{1.0cm}
{\Large\bfseries A PROJECT STAGE-1 REPORT\par}
\vspace{0.5cm}
{\large on\par}
\vspace{0.5cm}
{\fontsize{18}{22}\selectfont\bfseries \ProjectTitle\par}
\vspace{0.8cm}
{\large\bfseries By\par}
\vspace{0.3cm}
\begin{tabular}{>{\bfseries}l @{\hspace{1.5cm}} l}
\StudentOne & Roll No. \RollOne\\
\StudentTwo & Roll No. \RollTwo\\
\StudentThree & Roll No. \RollThree\\
\end{tabular}
\vspace{0.8cm}

{\large B.E. (ELECTRONICS \& TELECOMMUNICATION)\par}
\vspace{0.4cm}
{\large Under the guidance of\par}
{\Large\bfseries \GuideName\par}
\vspace{0.7cm}
% Uncomment and use the exact logo filename if available:
% \includegraphics[width=2.5cm]{clg-logo.png}\\[0.5cm]
{\large\bfseries DEPARTMENT OF ELECTRONICS \& TELECOMMUNICATION ENGINEERING\par}
{\large\bfseries \CollegeName\par}
{\large \CollegeAddress\par}
{\small Affiliated to Savitribai Phule Pune University, Pune\par}
\vfill
{\large\bfseries (\AcademicYear)\par}
\end{center}
\end{titlepage}

% =========================================================
% FORMAL TITLE PAGE
% =========================================================
\begin{titlepage}
\thispagestyle{empty}
\begin{center}
\vspace*{1cm}
{\fontsize{18}{22}\selectfont\bfseries \ProjectTitle\par}
\vspace{1cm}
A Project Report submitted in partial fulfillment of the requirements for award of the degree
\vspace{0.5cm}

{\Large\bfseries BACHELOR OF ENGINEERING\par}
\vspace{0.2cm}
{\large In\par}
{\Large\bfseries ELECTRONICS \& TELECOMMUNICATION\par}
{\large Of\par}
{\large\bfseries SAVITRIBAI PHULE PUNE UNIVERSITY\par}
\vspace{0.6cm}
{\large\bfseries By\par}
\vspace{0.3cm}
\begin{tabular}{>{\bfseries}l @{\hspace{1.5cm}} l}
\StudentOne & Roll No. \RollOne\\
\StudentTwo & Roll No. \RollTwo\\
\StudentThree & Roll No. \RollThree\\
\end{tabular}
\vspace{0.7cm}

Under the guidance of\\[0.2cm]
{\Large\bfseries \GuideName\par}
\vfill
% \includegraphics[width=2.5cm]{clg-logo.png}\\[0.4cm]
{\large\bfseries DEPARTMENT OF ELECTRONICS \& TELECOMMUNICATION ENGINEERING\par}
{\large\bfseries \CollegeName\par}
{\large \CollegeAddress\par}
\vspace{0.4cm}
{\large\bfseries (\AcademicYear)\par}
\end{center}
\end{titlepage}

% =========================================================
% CERTIFICATE
% =========================================================
\clearpage
\thispagestyle{empty}
\begin{center}
% \includegraphics[width=2.2cm]{clg-logo.png}\\[0.4cm]
{\Large\bfseries CERTIFICATE\par}
\vspace{0.8cm}
This is to certify that the Project Stage-1 entitled\\[0.4cm]
{\Large\bfseries ``\ProjectTitle''\par}
\vspace{0.5cm}
Submitted By\\[0.4cm]
\StudentOne\\
\StudentTwo\\
\StudentThree
\end{center}

\vspace{0.5cm}
has successfully completed the project Stage-1 under the supervision of
\textbf{\GuideName} in partial fulfillment of the requirements for the degree of
Bachelor of Engineering in Electronics \& Telecommunication of Savitribai Phule
Pune University. This work has not been submitted elsewhere for any degree.

\vspace{1.8cm}
\noindent
\begin{minipage}[t]{0.30\textwidth}\centering
\textbf{\GuideName}\\Guide
\end{minipage}\hfill
\begin{minipage}[t]{0.30\textwidth}\centering
\textbf{\CoordinatorName}\\Project Coordinator
\end{minipage}\hfill
\begin{minipage}[t]{0.30\textwidth}\centering
\textbf{\HODName}\\H.O.D. (E\&TC)
\end{minipage}

\vspace{2cm}
\noindent
\begin{minipage}[t]{0.45\textwidth}\centering Examiner\end{minipage}\hfill
\begin{minipage}[t]{0.45\textwidth}\centering
\textbf{\PrincipalName}\\Principal\\SCOE, Pune
\end{minipage}

% =========================================================
% FRONT MATTER
% =========================================================
\clearpage
\pagenumbering{roman}
\tableofcontents
\clearpage
\listoffigures
\addcontentsline{toc}{chapter}{List of Figures}
\clearpage

\chapter*{Abstract}
\addcontentsline{toc}{chapter}{Abstract}
The Intelligent Pre-Toll Processing and Dynamic Lane Management System is designed
to reduce toll-plaza congestion and vehicle waiting time through automated vehicle
verification. The system combines RFID-based vehicle identification with camera-based
number-plate recognition for dual verification. An ESP32 controller compares the RFID
and number-plate data to identify mismatched or unauthorized vehicles. An OLED
display provides local status information, while Blynk IoT can support remote
monitoring. The system also proposes dynamic lane management and automated
toll-barrier control using servo motors, providing an efficient and reliable
toll-management prototype.

\chapter*{Acknowledgement}
\addcontentsline{toc}{chapter}{Acknowledgement}
We are grateful to our project guide, \textbf{\GuideName}, for valuable guidance,
patience, and support throughout the project. We also express our sincere thanks to
\textbf{\CoordinatorName}, Project Coordinator, Department of Electronics \&
Telecommunication, for encouragement, suggestions, and coordination.

We thank \textbf{\HODName}, Head of the Department, and
\textbf{\PrincipalName}, Principal of \CollegeName, for their support and for
providing the facilities required to carry out this project. Finally, we thank our
families, friends, and everyone who encouraged us during the work.

\vspace{0.8cm}
\begin{flushright}
\StudentOne\\
\StudentTwo\\
\StudentThree
\end{flushright}

\chapter*{Abbreviations}
\addcontentsline{toc}{chapter}{Abbreviations}
\begin{description}[style=nextline]
\item[ANPR] Automatic Number Plate Recognition
\item[AI] Artificial Intelligence
\item[CNN] Convolutional Neural Network
\item[CPU] Central Processing Unit
\item[GPU] Graphics Processing Unit
\item[IoT] Internet of Things
\item[OCR] Optical Character Recognition
\item[OLED] Organic Light-Emitting Diode
\item[RFID] Radio Frequency Identification
\item[SPI] Serial Peripheral Interface
\item[UID] Unique Identifier
\item[Wi-Fi] Wireless Fidelity
\item[YOLO] You Only Look Once
\end{description}

\clearpage
\pagenumbering{arabic}

% =========================================================
% CHAPTER 1
% =========================================================
\chapter{Introduction}
\section{Introduction}
The Intelligent Pre-Toll Processing and Dynamic Lane Management System is a smart
transportation solution developed to improve the efficiency of toll-plaza operations.
Conventional toll plazas may experience congestion and increased waiting time due to
manual verification, vehicle queues, and delays in toll processing. The proposed
system aims to reduce these problems by introducing automated vehicle identification
and verification before a vehicle reaches the toll barrier.

The system uses an RFID reader to identify vehicles through their RFID tags. A camera
captures the vehicle's number plate, which can be processed using image-processing
and Automatic Number Plate Recognition (ANPR) techniques. The ESP32 microcontroller
acts as the embedded controller. The RFID-linked vehicle record is compared with the
number plate detected by the computer-vision software to perform dual verification.
This can help identify mismatched, unauthorized, or exceptional vehicles.

The system also proposes dynamic lane management to improve vehicle movement
through the toll plaza. Servo motors represent automated toll barriers, while an OLED
display provides vehicle and lane status. Wi-Fi and a platform such as Blynk may be
used for remote monitoring if implemented. The current prototype is an academic
demonstration; the final report should distinguish completed functions from planned
integration.

\section{Aim}
To design and develop an Intelligent Pre-Toll Processing and Dynamic Lane Management
System using ESP32, RFID, and camera-based number-plate recognition to provide
automated vehicle verification, reduce waiting time, and improve toll-plaza traffic
management.

\section{Objectives}
\begin{enumerate}
\item Identify vehicles using RFID-based identification.
\item Capture vehicle images and recognize number plates using computer vision and OCR.
\item Perform dual verification by comparing RFID-linked records with number-plate data.
\item Use an ESP32 microcontroller for vehicle-verification logic and system control.
\item Detect RFID--vehicle mismatches and reject unverified vehicles.
\item Demonstrate lane selection and management using two model lanes.
\item Control toll barriers automatically using servo motors.
\item Provide local status information through an OLED display.
\item Evaluate the prototype using valid and invalid test cases.
\end{enumerate}

\section{Block Diagram}
The proposed system includes an RFID tag and reader, a camera-based computer-vision
module, an ESP32 controller, an OLED display, and two servo-operated lane barriers.
The camera software sends the plate-verification result to the controller through a
communication interface that must be specified in the final implementation.

\begin{figure}[H]
\centering
\projectimage{block-diagram.png}{0.85}
\caption{Block diagram of the Intelligent Pre-Toll Processing and Dynamic Lane Management System.}
\label{fig:block}
\end{figure}

% =========================================================
% CHAPTER 2
% =========================================================
\chapter{Literature Review}
\section{Introduction}
The literature review covers technologies related to intelligent transportation,
RFID-based vehicle identification, wireless communication, and automatic number-plate
recognition. Existing research explores toll collection and traffic management using
RFID, computer vision, IoT, and artificial intelligence. Studying these systems helps
identify their operating principles, advantages, limitations, and possible improvements.

\section{Literature Review}
Automatic number-plate recognition systems commonly use object detection, image
processing, and OCR to locate a plate and extract its registration text. Recent object
detection approaches, including YOLO-based models, are designed to identify objects
in images or video frames efficiently. A dedicated plate-detection stage is still needed
when the selected model detects vehicles but does not localize number plates.

RFID technology provides contactless identification through a reader and tag. The
identifier can be mapped to a registered vehicle record. However, a tag identifier by
itself does not visually confirm that the vehicle passing through the lane matches the
record. Combining RFID with number-plate recognition adds a second verification step.

ESP32 development boards provide GPIO, SPI, I\textsuperscript{2}C, and wireless
communication capabilities suitable for embedded prototypes. In this project, the
ESP32 is used for RFID processing, display updates, and servo control, while the
computer performs computationally intensive image processing.

\section{Literature Review Summary}
\begin{longtable}{p{0.06\textwidth}p{0.30\textwidth}p{0.22\textwidth}p{0.12\textwidth}p{0.22\textwidth}}
\caption{Summary of related technologies and literature}\\
\toprule
\textbf{No.} & \textbf{Topic / Title} & \textbf{Author / Source} & \textbf{Year} & \textbf{Method / Principle}\\
\midrule
\endfirsthead
\toprule
\textbf{No.} & \textbf{Topic / Title} & \textbf{Author / Source} & \textbf{Year} & \textbf{Method / Principle}\\
\midrule
\endhead
1 & Automatic Number Plate Detection and Recognition Using YOLOv8 and CNN & Gajurel et al. & 2024 & Vehicle/plate detection and character recognition.\\
2 & ESP32 Wi-Fi microcontroller documentation & Espressif Systems & Check cited version & Wireless communication and embedded IoT functions.\\
3 & A Conceptual Study of Long-Range Active RFID System for Reliable Data Communication & Abdulla R. & 2014 & Long-range RFID identification and communication.\\
\bottomrule
\end{longtable}

\section{Summary}
The reviewed technologies indicate that RFID can provide contactless identification,
computer vision can detect vehicles and read number plates, and an ESP32 can control
the physical outputs of a prototype. The proposed project combines these functions
to demonstrate dual verification and model lane-barrier control. All final references
should be checked against the original publisher or official documentation before
submission.

% =========================================================
% CHAPTER 3
% =========================================================
\chapter{Methodology}
\section{Introduction}
The proposed system uses AI-based vehicle detection and number-plate recognition to
support automated vehicle identification. A camera captures images or video of vehicles
approaching the toll lane. The captured data is processed using computer-vision
techniques. The system detects the vehicle and then localizes its number plate. OCR
extracts the registration text. The extracted information is compared with the vehicle
record associated with the RFID tag.

\section{Methodology 1: Vehicle Detection}
The first stage focuses on vehicle detection using a suitable YOLO object-detection
model. A camera captures the road scene and provides images or video frames as input.
The model processes each frame and identifies vehicle classes such as cars, motorcycles,
buses, or trucks, depending on the model and its training classes. Bounding boxes and
confidence scores are produced for detections.

\subsection{Part 1: Vehicle Detection Implementation}
The vehicle-detection module is implemented in Python. Input images or camera frames
are processed by the selected model. For every detected vehicle, the program displays
a bounding box, class label, and confidence score. This stage provides vehicle
localization for subsequent processing. A vehicle detector does not automatically
guarantee number-plate detection; a plate-specific model or localization stage is
required separately.

\subsection{Algorithm}
The general detection algorithm follows these steps:
\begin{enumerate}
\item Capture an image or video frame from the camera.
\item Preprocess the frame if required by the selected model.
\item Pass the frame through the object-detection model.
\item Obtain predicted bounding boxes, class labels, and confidence scores.
\item Apply non-maximum suppression if supported by the detection pipeline.
\item Display or store the final vehicle detections.
\item Pass the relevant vehicle region to the plate-detection stage.
\end{enumerate}

\subsection{Flow Chart}
\begin{figure}[H]
\centering
\projectimage{flowchart.png}{0.70}
\caption{Flowchart of vehicle and number-plate detection.}
\label{fig:flowchart}
\end{figure}

\section{Methodology 2: Number-Plate Detection and Recognition}
After a vehicle is detected, the system identifies the number-plate region from the
captured image. Image processing may be used to crop, resize, and improve the plate
region. OCR is applied to recognize the characters present on the plate. The extracted
text is normalized, for example by converting letters to uppercase and removing spaces,
before comparison with the registration number linked to the RFID tag.

The system must reject an unreadable or low-confidence result instead of treating it
as a successful match. Recognition accuracy depends on image resolution, plate angle,
lighting, motion blur, plate condition, and the selected detector/OCR model.

\section{RFID Verification and Barrier Control}
The RFID reader obtains a tag UID and the controller searches the registered sample
records. If the UID is unknown, the transaction is rejected. If the UID is known, the
computer vision result is checked against the stored plate. When the plate matches and
the simulated balance is sufficient, the controller deducts the configured toll amount
and opens the selected servo barrier. Otherwise, the barrier remains closed and the
OLED displays the reason for rejection.

\section{System Algorithm}
\begin{enumerate}
\item Start the system and initialize RFID, OLED, servo, and communication interfaces.
\item Wait for an RFID tag and read its UID.
\item Check whether the UID is registered.
\item If the UID is unknown, display rejection and keep both barriers closed.
\item Obtain the camera/OCR result from the computer.
\item If the result is missing or unreliable, reject the transaction.
\item Compare the normalized plate text with the plate stored for the RFID UID.
\item If the plates do not match, reject the transaction.
\item Check the simulated account balance.
\item If the balance is sufficient, deduct the simulated toll and select a lane.
\item Open the selected servo barrier for the demonstration interval.
\item Close the barrier and return to the waiting state.
\end{enumerate}

% =========================================================
% CHAPTER 4
% =========================================================
\chapter{Results and Discussions}
\section{Results}
This chapter must be updated with results actually observed during testing. Insert
screenshots from the running vehicle-detection program, plate-detection output,
ESP32 serial monitor, and Wokwi or physical hardware simulation.

\begin{figure}[H]
\centering
\projectimage{plate-detection-result.png}{0.75}
\caption{Testing result of number-plate detection.}
\label{fig:plate-result}
\end{figure}

The number-plate detection module should be evaluated using different vehicle images
and input frames. Record the number of plates successfully detected and recognized,
along with examples of failures. Do not claim a specific recognition accuracy unless
it has been measured using a defined test set.

\begin{figure}[H]
\centering
\projectimage{vehicle-detection-result.png}{0.75}
\caption{Testing result of vehicle detection.}
\label{fig:vehicle-result}
\end{figure}

The vehicle-detection module can be tested with images containing different vehicle
types. Record which classes are detected, whether the bounding boxes are correctly
positioned, and how lighting or image quality affects performance.

\section{Simulation}
\begin{figure}[H]
\centering
\projectimage{wokwi-simulation.png}{0.75}
\caption{Wokwi simulation of the ESP32-based toll controller.}
\label{fig:simulation}
\end{figure}

\subsection{Simulation Results}
\begin{figure}[H]
\centering
\projectimage{simulation-result.png}{0.75}
\caption{Simulation result and controller status.}
\label{fig:simulation-result}
\end{figure}

\begin{table}[H]
\centering
\caption{Functional test cases and expected behaviour}
\begin{tabular}{p{0.08\textwidth}p{0.42\textwidth}p{0.38\textwidth}}
\toprule
\textbf{ID} & \textbf{Test condition} & \textbf{Expected result}\\
\midrule
T1 & Registered RFID and matching plate & Approve transaction and open selected barrier\\
T2 & Unknown RFID UID & Reject; barrier remains closed\\
T3 & Registered RFID with mismatched plate & Reject; barrier remains closed\\
T4 & Camera/OCR result unavailable or unreliable & Reject; barrier remains closed\\
T5 & Insufficient simulated balance & Reject; do not deduct toll or open barrier\\
T6 & Successful lane operation & Selected servo opens and closes\\
T7 & OLED status display & Appropriate transaction status appears\\
\bottomrule
\end{tabular}
\end{table}

\section{Discussions}
The proposed system combines RFID identification with an independent visual plate
check. The ESP32 provides embedded control and local status feedback, while a computer
performs the image-processing tasks. The reliability of the overall system depends on
the quality of plate detection and OCR and on successful communication between the
computer and ESP32.

Potential limitations include poor lighting, glare, angled or dirty plates, motion blur,
OCR errors, duplicate RFID reads, and communication failures. The prototype should
fail closed when the verification result is uncertain. Real-world deployment would
require secure tag authentication, secure records, transaction logging, privacy
protection, reliable vehicle-passage sensing, and regulatory approval.

% =========================================================
% CHAPTER 5
% =========================================================
\chapter{Conclusions}
\section{Conclusions}
The Intelligent Pre-Toll Processing and Dynamic Lane Management System proposes a
combined approach to vehicle identification using RFID and camera-based number-plate
recognition. The computer-vision module is intended to detect vehicles and read their
registration plates, while the ESP32 handles RFID record lookup, status display, simulated
toll deduction, and servo-barrier control.

The main contribution of the prototype is the additional verification rule that the
number plate detected by the camera must match the registration number associated
with the RFID tag. Unknown tags, mismatched plates, and failed camera verification
are rejected. The final conclusion must be based on the integrated prototype and the
test results recorded by the project team.

\section{Future Scope}
The system can be enhanced in the following ways:
\begin{itemize}
\item Improve number-plate detection and OCR using a larger and more diverse dataset.
\item Integrate the camera software and ESP32 using a defined communication protocol.
\item Add a database for registered vehicles, simulated balances, and transaction logs.
\item Implement queue measurement and a documented dynamic lane-selection algorithm.
\item Add vehicle-passage sensors to close barriers based on actual passage rather than a fixed delay.
\item Add secure authentication, fault recovery, and privacy controls.
\item Evaluate performance under different lighting, camera angles, vehicle speeds, and plate conditions.
\item Investigate integration with FASTag or authorized toll databases only where suitable access and approvals exist.
\end{itemize}

% =========================================================
% REFERENCES
% =========================================================
\clearpage
\begin{thebibliography}{99}
\addcontentsline{toc}{chapter}{References}

\bibitem{gajurel}
K. Gajurel, M. K. Baniya, K. Mahato, M. Kathet, S. Lingden, and P. Karki,
``Automatic Number Plate Detection and Recognition Using YOLOv8 and CNN,''
\textit{Proceedings of the 15th IOE Graduate Conference}, vol. 15, pp. 143--149, 2024.
Verify the final bibliographic details against the original publication.

\bibitem{esp32}
Espressif Systems, \textit{ESP32 Series Datasheet} and official ESP32 documentation.
Use the documentation for the exact development board selected.

\bibitem{abdulla}
R. Abdulla, ``A Conceptual Study of Long Range Active RFID System for Reliable
Data Communication,'' in \textit{Proceedings of the International Conference on
Frontiers of Communications, Networks and Applications (ICFCNA)}, 2014, pp. 1--6.
Verify the publication details before submission.

\bibitem{opencv}
OpenCV, \textit{Open Source Computer Vision Library Documentation}.
Available at: \url{https://docs.opencv.org/}.

\bibitem{tesseract}
Tesseract OCR, \textit{Tesseract User Manual}.
Available at: \url{https://tesseract-ocr.github.io/}.

\bibitem{ultralytics}
Ultralytics, \textit{YOLO Documentation}.
Available at: \url{https://docs.ultralytics.com/}. Cite the exact model and version used.

\bibitem{mfRC522}
NXP Semiconductors, \textit{MFRC522 Standard Performance MIFARE and NTAG Frontend:
Datasheet}. Use the manufacturer's datasheet for the exact module specifications.

\end{thebibliography}

\end{document}
